\section{Summary}

This qualification document presented a doctoral research proposal centered on
source code minimaps as visual signatures for software analysis. The proposal
starts from a practical problem: modern software repositories are heterogeneous,
large, and difficult to understand at scale, while many analysis approaches
depend on language-specific processing, expert-defined rules, or representations
that preserve sensitive lexical content.

The approach rests on the commitments stated in
Chapter~\ref{chapter:introduction}: each source file is rendered as a
two-dimensional image whose geometry comes from the layout developers actually
write; the image is produced for a model to read, not for a human to inspect;
and the lexical layer that carries business-sensitive content is suppressed
before anything leaves the developer's environment. Informal evidence discussed
in Chapter~\ref{chapter:preliminary-results} suggests that this suppression
costs little predictive performance, although the controlled ablation that will
quantify this cost is still pending.

The resulting images can be analyzed with handcrafted texture descriptors, deep
learning models, and explainable AI methods. The approach does not replace
static analysis, parsers, or software visualization; it complements them in
contexts where sending raw source code to external services is not acceptable.

\section{Current Evidence}

The research has already produced preliminary evidence. The first empirical
study showed that Local Binary Patterns and traditional classifiers can detect
software features from source code minimaps, with Random Forest reaching
approximately 93\% F1-score in the evaluated setting. The second study scaled
the representation to 685,906 minimaps from 100 open-source projects across 10
programming languages and evaluated deep models for Type, Project, and Author
objectives. The tool-support line, MinimapSignatures, demonstrated that minimap
generation and inference can be integrated into Visual Studio Code and IntelliJ
IDEA through a shared backend.

These results support the central hypothesis that anonymized minimaps preserve
useful structural information for software analysis. They also reveal the main
remaining challenges: class imbalance, label quality, anonymization strength,
explainability, generalization, and developer-centered evaluation.

\section{Expected Thesis Contribution}

The expected thesis contribution is a coherent methodology and artifact set for
visual source code analysis with minimaps: a representation, an anonymization
strategy, dataset-generation procedures, empirical evaluation, explainability
protocols, and IDE-integrated tooling. Scientifically, the thesis aims to
position source code minimaps as a reproducible bridge between software
visualization and machine learning for code. Practically, it aims to provide
tools and datasets that can be reused in future research on software
comprehension, quality assessment, architectural deviation, and code-review
support.

The qualification stage shows that the project has moved beyond an initial
idea: it already includes empirical studies, a large dataset, explainability
results, and working tool prototypes. The remaining doctoral work will refine
the evidence, expand the task set, evaluate practical usefulness, and
consolidate the approach into a thesis. By showing that structural information
alone supports a family of software analysis tasks, the thesis will offer both
a practical option for confidentiality-aware analysis and a scientific argument
that layout can stand in for lexical content in this class of problems.
